\section{Background and Related Work}\label{sec:background}

This section defines the SQL antipatterns evaluated in the study and explains how jOOQ represents them in Java source code. It then compares the source representations and evaluation units used by prior detection approaches to identify the occurrence-localisation gap addressed by this study.

% Source: chapters/02_background.tex, SQL Antipatterns; reduced and operationalised in Step 6.
\subsection{SQL antipatterns and evaluated classes}\label{sec:sqlAntipatterns}

SQL antipatterns are recurring database-design or query-construction choices that impair properties such as performance, maintainability, portability, or data integrity \citetext{\citealp[p. 1]{karwin2023sql}; \citealp[p. 4]{alshemaimri2021survey}}. Karwin's catalogue, first published in 2010, provides the base definitions used in this study \nobreakcite{karwin2023sql}. Later work has proposed broader and vendor-specific taxonomies \nobreakcite{alshemaimri2021survey,10.1007/978-3-031-09070-7_26,10.1007/978-3-031-35311-6_51,factor2014sql,angelakos2025PostgreSQL}.

Prior studies found SQL antipatterns in open-source and industrial systems \nobreakcite{muse2020prevalence,sharma2018smelly}. In SQLite-based Android applications, removing Implicit Columns reduced mean query runtime by \qty{27}{\percent} and energy use by \qty{29}{\percent} \nobreakcite[p. 60]{lyu2019quantifying}.

The detector experiments retain seven classes. Table~\ref{tab:evaluatedAntipatterns} states this study's operational definitions, including the boundaries that differ from the broader conceptual definitions. The full Karwin taxonomy and the class-specific annotation and prompt decision rules are available in the online replication materials.

\begin{longtable}{@{}p{0.24\textwidth}p{0.70\textwidth}@{}}
    \caption{Operational definitions of the seven SQL antipattern classes retained for detector evaluation.}
    \label{tab:evaluatedAntipatterns}\\
    \hline
    \thead[l]{Class} & \thead[l]{Operational definition} \\ \hline
    \endfirsthead
    \caption[]{Operational definitions of the seven retained classes (continued).}\\
    \hline
    \thead[l]{Class} & \thead[l]{Operational definition} \\ \hline
    \endhead
    \hline
    \endfoot
    \hline
    \endlastfoot
    ID Required &
    A generated table class, excluding views, defines a primary key named exactly \texttt{id}, ignoring case, or defines a synthetic primary key despite a suitable natural or compound unique key \nobreakcite[Ch. 4]{karwin2023sql}. \\

    Keyless Entry &
    A generated table class, excluding views, contains a column interpreted as referring to another table but lacks a foreign-key constraint, provided the supplied key definitions contain a suitable referenced primary key \nobreakcite[Ch. 5]{karwin2023sql}. \\

    Rounding Errors &
    A schema stores an exact fractional quantity in \texttt{FLOAT} or \texttt{REAL} rather than \texttt{DECIMAL} or \texttt{NUMERIC} \nobreakcite[Ch. 10]{karwin2023sql}. \\

    31 Flavors &
    A column encodes a fixed allowed-value list with a \texttt{CHECK} constraint or \texttt{ENUM} instead of a lookup table. Emptiness checks and range constraints are excluded \nobreakcite[Ch. 11]{karwin2023sql}. \\

    Fear of the Unknown &
    A schema uses a semantically meaningless sentinel for missing data, marks a practically non-null column as nullable, or jOOQ query logic applies null-unsafe logic to a nullable database column. Java-only null handling and query values supplied from Java are excluded \nobreakcite[Ch. 14]{karwin2023sql}. \\

    Poor Man's Search Engine &
    A jOOQ condition uses \texttt{LIKE}, \texttt{ILIKE}, or a regular expression for full-text matching. Prefix searches are excluded. Calls whose parameters do not reveal whether full-text wildcards are present are included by the detector rule \nobreakcite[Ch. 17]{karwin2023sql}. \\

    Implicit Columns &
    A jOOQ query or generated-DAO operation fetches all columns without an explicit projection. Calls within \texttt{fetchCount} or \texttt{fetchExists} are excluded \nobreakcite[Ch. 19]{karwin2023sql,jooq2025select}. \\
\end{longtable}

% Source: chapters/02_background.tex, jOOQ Object Oriented Querying; reduced in Step 6.
\subsection{jOOQ source representation}\label{sec:jooq}

jOOQ represents SQL as a Java domain-specific language \nobreakcite{jooq2025sql}. Its code generator maps database objects, including tables and constraints, to Java classes \nobreakcite{jooq2025codegen}. Projects may also execute plain SQL through its JDBC wrapper \nobreakcite{jooq2025execution}.

Implicit Columns denotes a blind projection, such as \texttt{SELECT *}, whose result can change when a table's columns change \nobreakcite[Ch. 19]{karwin2023sql}. Figure~\ref{fig:implicitColumnsRepresentations} shows the same blind projection as SQL text and as a jOOQ DSL expression.

\begin{figure}[h]
    \textbf{SQL}
    \begin{lstlisting}[frame=none, basicstyle=\small]
SELECT * FROM employee WHERE name = 'Jack'
    \end{lstlisting}

    \textbf{jOOQ}
    \begin{lstlisting}[frame=none, basicstyle=\small]
var employee = dslContext.select(DSL.asterisk())
    .from(EMPLOYEE)
    .where(EMPLOYEE.NAME.eq("Jack"))
    .fetchOne();
    \end{lstlisting}
    \caption{Equivalent Implicit Columns occurrences in SQL text and a jOOQ Java DSL expression. Both select every column from the \texttt{employee} table.}
    \label{fig:implicitColumnsRepresentations}
\end{figure}

A jOOQ statement is assembled from Java calls and may refer to schema information through generated classes. The corresponding SQL is produced at runtime \nobreakcite{jooq2025withgen}. Source-level detection must therefore interpret DSL calls and generated schema representations, rather than relying only on recoverable SQL strings.

% Source: chapters/02_background.tex, Static Analysis for Antipattern Detection and Detection of SQL Antipatterns; synthesised in Step 6.
\subsection{Detection approaches and the localisation gap}\label{sec:staticAnalysis}

Static detectors commonly encode metric thresholds or rules over source representations such as abstract syntax trees (ASTs) \nobreakcite{DBLP:conf/icsm/Marinescu04,DBLP:journals/tse/MohaGDM10}. Metric and pattern rules can produce false positives when the required context is absent \nobreakcite[pp. 10--11]{DBLP:journals/jserd/PaivaDFS17}. Learned code-smell detectors instead infer patterns from annotated code \nobreakcite{DBLP:journals/ese/FontanaMZM16}.

Two SQL-antipattern tools illustrate extraction-and-parsing approaches for SQL embedded in host-language source. DbDeo extracts SQL strings, parses them, and applies schema-smell rules \nobreakcite{sharma2018smelly}. SQLInspect uses intra-procedural string resolution and a fault-tolerant parser before applying query-antipattern rules \nobreakcite{nagy2017static,nagy2018sqlinspect}. SQLCheck analyses application query logs with static and dynamic checks \nobreakcite{DBLP:conf/sigmod/DintyalaNA20}. These tools expect SQL text or executed queries. Their coverage of source-level jOOQ DSL expressions and generated schema classes has not been established.

Recent studies apply LLMs to code-quality detection \nobreakcite{blaauwendraad2025evaluating}. Most directly, \nobreaktextcite[p. 410]{de2025effectiveness} evaluated language models for PL/SQL smell detection. Table~\ref{tab:detectionApproaches} compares the representations and output units of these approaches with the requirements of this study.

\begin{longtable}{@{}p{0.20\textwidth}p{0.34\textwidth}p{0.40\textwidth}@{}}
    \caption{Detection representations and output units relevant to jOOQ source analysis.}
    \label{tab:detectionApproaches}\\
    \hline
    \thead[l]{Approach} & \thead[l]{Representation and output} & \thead[l]{Relevance to this study} \\ \hline
    \endfirsthead
    \caption[]{Detection representations and output units relevant to jOOQ source analysis (continued).}\\
    \hline
    \thead[l]{Approach} & \thead[l]{Representation and output} & \thead[l]{Relevance to this study} \\ \hline
    \endhead
    \hline
    \endfoot
    \hline
    \endlastfoot
    SQL extraction and parsing &
    SQL recovered from host-language source or query logs; parsed-statement warnings \nobreakcite{sharma2018smelly,nagy2017static,DBLP:conf/sigmod/DintyalaNA20}. &
    Does not directly expose occurrences represented by jOOQ DSL calls or generated schema classes. \\

    API or AST rules &
    Host-language syntax and API patterns; warnings attached to matched source nodes \nobreakcite{DBLP:journals/tse/MohaGDM10}. &
    Can localise encoded patterns precisely, but requires explicit coverage of jOOQ idioms and cross-file relationships. No evaluated seven-class jOOQ baseline was found in the reviewed work. \\

    LLM source analysis &
    Source code and supplied context; prior work reports smell recognition or classification \nobreakcite{de2025effectiveness,blaauwendraad2025evaluating}. &
    Can return class-labelled spans, but span agreement and class-specific errors require empirical evaluation. \\

    This study &
    Java DSL and generated schema code; multiple class-labelled source-line spans. &
    Evaluates the same occurrence unit later counted as detector flags across repositories. \\
\end{longtable}

A file may contain several occurrences of the same antipattern. File-level classification establishes only whether a class is present; it does not establish how many distinct occurrences the detector recovers or whether their source spans agree with the reference annotations. The reviewed work does not report occurrence-level localisation for the seven retained classes in jOOQ code. This study therefore evaluates same-file, same-class line-span agreement before using detector flags in the repository analysis.
